%% Contents -- original pages ix-xii (serif, added in the 1998 edition)
\clearpage
\def\chapcode{}%
\renewcommand\capsshorttitle{Contents}%
\thispagestyle{fmplain}%
\pagestyle{fmfancy}%
\pdfbookmark[0]{Contents}{contents}%
\phantomsection
\begin{center}
  \rule{0.85\textwidth}{0.5pt}\\[1.5pt]
  \rule{0.85\textwidth}{2pt}\\[1.2\baselineskip]
  {\Huge\bfseries Contents}
\end{center}
\vspace{3\baselineskip}
\makeatletter
%% \tableofcontents uses exactly \@starttoc{toc}: it first reads the
%% previous run's \jobname.toc (so entries typeset), then opens that file
%% for *writing*, so this run's \addtocontents lines land in main.toc for
%% the next pass.  A plain \@input@ would read-only and never populate the
%% file, which is why main.toc did not exist before.
\@starttoc{toc}%
\makeatother
\clearpage
